%------------------------------------------------------------------------------%
\documentclass[a4paper,12pt,fleqn]{article}

\usepackage{gaal}

\ShowAnswers

%------------------------------------------------------------------------------%
\begin{document}

\makeHeader{GAAL}{Prova 3 -- Turma 4}{26/10/2026}

% 01
%------------------------------------------------------------------------------%
\exercise{20\%}
Determine um vetor normal e uma equação do plano que passa pelos pontos
\(P = (3, 3, 1)\),
\(Q = (0, 3, 2)\) e
\(R = (1, 0, 1)\)

\begin{answer}
  Calculamos dois vetores contidos no plano
  \[
    \overrightarrow{PQ}
    = Q - P
    = (0, 3, 2) - (3, 3, 1)
    =( -3, 0, 1)
  \]
  \[
    \overrightarrow{PR}
    = R - P
    = (1, 0, 1) - (3, 3, 1)
    = (-2, -3, 0)
  \]

  Um vetor normal ao plano é dado pelo produto vetorial
  \begin{align*}
    \vec{n}
    & = \overrightarrow{PQ}\times\overrightarrow{PR}
      = \begin{vmatrix}
      \vec{i} & \vec{j} & \vec{k} \\
      -3      &  0      & 1       \\
      -2      & -3      & 0
    \end{vmatrix}
    \\
    & = \big[  0\cdot 0 - 1\cdot(-3) \big] \vec{i}
      - \big[ -3\cdot 0 - 1\cdot(-2) \big] \vec{j}
      + \big[ -3(-3)    - 0(-2)      \big] \vec{k}
    \\
    & = (3, -2, -9)
  \end{align*}

  A equação do plano que passa por \(P = ( 3, 3, 1)\)
  e tem vetor normal \(\vec{n} = (3, -2, -9)\) é
  \begin{align*}
    \vec{n} \cdot (X-P) & = 0 \\
    (3, -2, -9) \cdot (x-3, y-3, z-1) & = 0 \\
    3(x-3) -2(y-3) -9 (z-1) & = 0 \\
    3x -9 -2y -6 -9z +9 & = 0 \\
    3x -2y -9z & = 6
    \end{align*}
\end{answer}

% 02
%------------------------------------------------------------------------------%
\exercise{20\%}
Determine a posição relativa entre os planos e uma equação paramétrica da interseção
\[
  \pi_1\colon  x + y + z = 3 \qquad
  \pi_2\colon 2x - y + z = 1
\]

\begin{answer}
  Os vetores normais $(1, 1, 1)$ e $(2, -1, 1)$ não são proporcionais.

  Os planos se intersectam em uma reta. Fazendo $z=t$:
  \begin{align*}
    x+y  & =3-t                      \\
    2x-y & =1-t                      \\
    3x   & =4-2t                     \\
    x    & =\frac{4}{3}-\frac{2}{3}t \\
    y    & =\frac{5}{3}-\frac{1}{3}t
  \end{align*}
  Assim,
  \[
    (x, y, z)
    = \left(\frac{4}{3},\frac{5}{3},0\right)+t\left(-\frac{2}{3},-\frac{1}{3},1\right)
  \]
\end{answer}

% 03
%------------------------------------------------------------------------------%
\exercise{20\%}
  Determine o ângulo entre os vetores
  $\vec{v} = (-1, 3, 2)$ e
  $\vec{w} = (2, -2, 4)$.

\begin{answer}
  \[
    \vec{v}\cdot\vec{w}
    =
  \]
  \[
    \norm{\vec{v}}
    = \norm{\vec{w}}
    =
  \]
  Logo
  \[
    \cos\theta
    =
  \]
  \[
    \theta = \arccos()
  \]
\end{answer}

% 04
%------------------------------------------------------------------------------%
\exercise{20\%}
Verifique se os vetores $\vec{v}_1=(1,0,1)$, $\vec{v}_2=(0,1,1)$ e $\vec{w}=(2,3,5)$
são linearmente independentes.

\begin{answer}
  Procuramos $a,b$ tais que
  \[
  a(1,0,1)+b(0,1,1)=(2,3,5)
  \]
  Comparando coordenadas, $a=2$, $b=3$ e $a+b=5$, que são compatíveis. Portanto,
  \[
  \vec{w}=2\vec{v}_1+3\vec{v}_2
  \]
  Para a independência, suponha $a\vec{v}_1+b\vec{v}_2=\vec{0}$.

  Então $(a,b,a+b)=(0,0,0)$, logo $a=b=0$.

  Os vetores são linearmente independentes.
\end{answer}

% 05
%------------------------------------------------------------------------------%
\exercise{20\%}
Obtenha, por Gram--Schmidt, uma base ortonormal que gere o mesmo espaço
que os vetores
$\vec{v}_1 = (1, 0, 1)$ e
$\vec{v}_2 = (1, 0, 0)$

\begin{answer}
  Tome
  \[
    \vec{u}_1=\vec{v}_1=(1,1)
    \qquad
    \vec{e}_1=\frac{1}{\sqrt{2}}(1,1)
  \]
  Agora,
  \begin{align*}
  \vec{u}_2&=\vec{v}_2-\frac{\vec{v}_2\cdot\vec{u}_1}{\norm{\vec{u}_1}^{2}}\vec{u}_1\\
  &=(1,0)-\frac{1}{2}(1,1)\\
  &=\left(\frac{1}{2},-\frac{1}{2}\right)
  \end{align*}
  Como $\norm{\vec{u}_2}=1/\sqrt{2}$,
  \[
  \vec{e}_2=\frac{1}{\sqrt{2}}(1,-1)
  \]
  A base ortonormal é $\{\vec{e}_1,\vec{e}_2\}$.

\end{answer}

%------------------------------------------------------------------------------%
\end{document}
%------------------------------------------------------------------------------%

